\documentclass[12pt, journal]{IEEEtran}
\usepackage{lipsum}
\usepackage{cite}
\usepackage{amsmath}
\usepackage{algorithmic}
\usepackage{array}
\usepackage{hyperref}
\usepackage{orcidlink}
\usepackage{stfloats}

\bibliographystyle{IEEEtran}

\usepackage[caption=false,font=footnotesize]{subfig}

\begin{document}

\markboth{Análisis de Señales Mixtas, IS 2026}{}

\title{Análisis de Señales Mixtas \\ \textit{Diseño de un sistema de compresión espectral adaptativa en tiempo real} }


\author{
\IEEEauthorblockA{\\Área Académica de Ingeniería en Computadores}
\IEEEauthorblockA{\\Instituto Tecnológico de Costa Rica} \\ 
\vspace{\baselineskip}\\
\IEEEauthorblockN{
Camila Lizano Brenes, 2024 \\
Dylan Guerrero, 20 \\
Javier Hernandez, 20 \\
Sergio Daniel Alvarez Chanto, 2024134332 \\}
}

\maketitle

\begin{abstract}
This paper presents the design and implementation of a real-time adaptive spectral compression system based on the Fast Fourier Transform (FFT). The system operates across three microcontrollers: a transmitter that captures a live analog signal via ADC, computes its FFT, and applies adaptive spectral compression by retaining only the highest-energy spectral coefficients; and two receivers that reconstruct the signal through the Inverse FFT (IFFT), one reproducing audio output, and the other computing quantitative quality metrics. The compression strategy leverages Parseval's theorem to identify the minimum number of spectral components required to preserve at least 95\% of the total signal energy. The system operates on power-of-two block sizes, consistent with the radix-2 Cooley-Tukey algorithm. Additionally, the relationship between frequency-domain multiplication and time-domain convolution is analyzed in the context of Linear Time-Invariant (LTI) systems. Results demonstrate that a 95\% energy threshold is highly effective for tonal signals but insufficient for broadband signals, showing that energy preservation percentage does not directly translate to perceptual reconstruction quality, with direct implications for embedded DSP design and low-bandwidth audio transmission.
\end{abstract}
 
\begin{IEEEkeywords}
Compresor, energía, transformada de Fourier, series de Fourier.
\end{IEEEkeywords}

\section{Introduction}
En este proyecto se busca diseñar un sistema de compresión espectral adaptativo en tiempo real. Para lograr esto, se debe generar una compresión y reconstrucción espectral en tiempo real utilizando la Transformada Rápida de Fourier (FFT) y la Transformada Inversa Rápida de Fourier (IFFT), respectivamente. El sistema completo cuenta con tres microcontroladores, cada uno con una funcionalidad independiente. El primero funciona como un transmisor, este es el encargado de captar una señal analógica (tono seno o señal compuesta) con un dispositivo ADC, calcular la FFT y realizar la compresion espectral adaptativa, la cual conserva solamente las componentes de mayor energía. El segundo y el tercer microcontrolador son receptores de la señal comprimida que transmite el primer microcontrolador. El segundo, al recibir la señal, realiza la reconstrucción en el dominio del tiempo aplicando la IFFT y estando conectado a un parlante DAC la reproduce, para verificar la similitud con la original. Finalmente, el tercer microcontrolador se encarga de reconstruir la señal, con el objetivo de calcular métricas cuántitativas como el Error Cuadrático Medio, la energía preservada y el error de la señal.\\
Al utilizar con conceptos matemáticos complejos para el desarrollo de este proyecto, se trabaja primero en una etapa investigativa, donde se analizan, comprenden y, finalmente, se aplican en la implementación del sistema. Este proceso se demuestra en este paper con la secuencia de apartados. Primero se presenta el marco teórico utilizado. Esto implica la literatura e información consultada y la explicación de los conceptos. A continuación, se presentarán los resultados y análisis obtenidos. Posteriormente, se presentarán conclusiones del desarrollo. Por último, se discutirán posibles mejoras.

\section{Marco teórico}
Como se menciona en la introducción, para el desarrollo de este proyecto se utilizaron múltiples conceptos matemáticos vitales a la hora de trabajar con señales analógicas no constantes. Se investigaron los siguientes conceptos: Series de Fourier, ortogonalidad, transformada discreta de Fourier, transformada rápida de Fourier y relación entre convolución y multiplicación en frecuencia. \\
La ortogonalidad es:\\
Las Series de Fourier son:\\
La Transformada de Fourier es:\\
La Transformada Discreta de Fourier es:\\
La Transformada rápida de Fourier es:\\
La relación entre convolución y multiplicación en frecuencia se presenta:\\

\section{Resultados y análisis}
A continuación se presentan los resultados obtenidos en la implementación del sistema previamente descrito.
Resultados de la implementación de la FFT y IFFT.\\
Por otro lado, los resultados obtenidos en la implementación del compresor de una señal fueron: 
(poner imagenes)

\section{Conclusiones}
En conclusión, en el apartado de la implementación de la FFT y IFFT se logró...\\
Además, en el apartado del compresor se logró comprimir una señal de audio captada por un dispositivo ADC y, esta misma, ser recontruida con un 95\% de la energía con un número menor de componentes. Lo cual generó, en la mayoría de veces, una pérdida en la calidad pero se mantenía percibible que era la señal original reconstruida.


\section{Mejoras}
posibles mejoras en prosa


\bibliography{references}

\end{document}
